\section{What the execution traces show}
\label{sec:mechanism}

% Remaster, 6 October 2026.

Harbor records every trial's execution trace. We check the traces for three
things: who builds the model and whether the discipline instruction changes
the agent's behaviour with the model builder unchanged (H2),
and how often an instruction sends the agent into repair (H3).

\subsection{Who builds the model: the agent or an AutoML library}
\label{sec:spontaneous}

Without a library instruction, both agents build the model themselves.
Terminus-2 fitted a gradient-boosting model on every TabReD task in every
trial: LightGBM in 23 of the 24 trials of Section~\ref{sec:abl-image} and
XGBoost in the remaining one. On the stock image it installed LightGBM
itself in 18 of the 24 trials. AutoDS-Tools without the library
instruction wrote LightGBM, XGBoost or CatBoost by hand in all 24 trials of
each of its two cells without it. The configuration is hand-written and
near-identical across tasks and trials: for Terminus-2 the code sets
1\,000 estimators and a learning rate of $0.05$ in all 24 trials and early
stopping on a holdout in 23. No trace of either agent contains a hyperparameter search,
by Optuna, grid, randomised search or a cross-validated sweep.

With the library instruction, AutoDS-Tools calls LightAutoML in every
trial that left a trace: 24 of 24 with the library instruction alone and 20
of 20 with the whole file. Four further trials with the whole file were cut
at the time limit before a trace was written. LightAutoML at its defaults fits a
linear model, LightGBM and CatBoost, tunes the two boosting models with
Optuna within its time budget and blends the results. A rerun of the
prescribed call on F-DATA recorded 101 Optuna trials per boosting model. The
library instruction therefore replaces one hand-written model with an AutoML
search. The strongest published method is tuned with Optuna and averaged
over fifteen seeds, and no agent ran a search of that size over its own
code. The per-trial extraction of learners and
settings is in the repository (\texttt{trace\_learners.csv}).

This is the mechanism behind \ref{hyp:builder}. Without a library the agent
writes one default configuration and stops. The library instruction moves
the search and the tuning into the model builder.


\subsection{The discipline instruction reaches the agent and leaves the score unchanged}
\label{sec:discipline}

The discipline instruction barely moves the task metric
(Table~\ref{tab:grid}), and the traces show that the agent follows it. Three clauses ask it to compare against a trivial predictor, to plan the
time budget before starting, and to keep the variant that validates best.
All three behaviours appear in every trial that left a trace in the two
cells with the discipline instruction. In the two cells without it, the
first two appear in no trial and the third in 2 of 48. The agent follows the
instruction, and the score stays where it was. This supports
\ref{hyp:builder} from the other side: knowledge that changes how the agent
works and leaves the model builder as it was leaves accuracy as it was.

\subsection{Instructions and execution failures}
\label{sec:friction}

AutoDS-Tools contains a repair loop, invoked when execution fails. The loop
was not logged in the repeated runs, so the count comes from an earlier pair
of runs. They cover all eighteen TabReD and MLAgentBench tasks once per
cell, on the earlier data preparation of Section~\ref{sec:coverage}. We
report only the direction. The loop fired 18 times with the file and 12 times without
it, mostly on tasks for which the file names a specialised library: four
invocations against none on \texttt{clrs} and three against none on
\texttt{amp-parkinsons}.

The largest loss on MLAgentBench traces back to one clause. The discipline
instruction tells the agent not to stop after ``an arbitrary small budget (5
or 15 epochs is almost always badly undertrained)'' and offers
``$\sim 100$ epochs'' as a starting point. On \texttt{clrs} the agent
without the file chose 5 to 30 epochs and used 92 of a 120-minute budget on
one trial. With the file it attempted a hundred epochs and made no
submission in three trials. For \ref{hyp:interaction}, an instruction can
therefore change reliability through execution failures. In AutoDS-Tools the
repair loop turns these failures into extra iterations, which is one reason
the file costs it few submissions. Where the loop cannot help, as on
\texttt{clrs} with a hundred epochs in a fixed budget, the trial is lost.
